\section{Prerequisites}
\label{sec:prerequisites}

This section fixes all notions used later. It is written bottom-up: no concept
is used before it is defined.

\subsection{Simplicial sets, horns, and quasicategories}

\begin{definition}[Simplicial set]
Let $\Delta$ denote the category whose objects are the finite nonempty linear
orders $[n] = \{0 < 1 < \dots < n\}$ and whose morphisms are order-preserving
maps. A \emph{simplicial set} is a functor $X : \Delta^{\mathrm{op}} \to
\mathsf{Set}$. The set $X_n = X([n])$ is the set of \emph{$n$-simplices};
$X_0$ are \emph{vertices}, $X_1$ are (directed) \emph{edges}. The category of
simplicial sets is denoted $\sset$.
\end{definition}

\begin{example}
The \emph{standard $n$-simplex} $\Delta^n = \Delta(-,[n])$. Concretely,
$\Delta^2$ has three vertices $0,1,2$, three edges $0\to1$, $1\to2$,
$0\to2$, and one nondegenerate $2$-simplex filling the triangle.
\end{example}

\begin{definition}[Horn]
For $0 \le k \le n$, the \emph{horn} $\Lambda^n_k \subset \Delta^n$ is the
subcomplex obtained by removing the interior of $\Delta^n$ and the face
opposite vertex $k$. A horn is \emph{inner} if $0 < k < n$.
\end{definition}

\begin{example}
$\Lambda^2_1$ consists of two composable edges $0 \to 1 \to 2$ with the
composite $0 \to 2$ and the filling triangle \emph{missing}. Filling
$\Lambda^2_1$ means supplying both.
\end{example}

\begin{definition}[Kan complex; quasicategory; nerve]
A simplicial set $X$ is a \emph{Kan complex} if every map $\Lambda^n_k \to X$
extends along $\Lambda^n_k \hookrightarrow \Delta^n$ for all horns, and a
\emph{quasicategory} if this holds for all \emph{inner} horns. The
\emph{nerve} $\N(\mathcal{C})$ of a small category $\mathcal{C}$ is the
simplicial set with $\N(\mathcal{C})_n = $ chains of $n$ composable
morphisms. A simplicial set is a nerve of a category iff every inner horn has
a \emph{unique} filler \citep{lurie2009htt}.
\end{definition}

\begin{remark}[The spectrum used in this paper]
Bare directed graphs (only $0$- and $1$-cells): nothing filled.
Quasicategories: everything filled, fillers unique only up to higher cells.
Nerves: everything filled uniquely. We will interpret these respectively as
the newborn universe, the solved universe, and the classical
(zero-information) limit.
\end{remark}

\subsection{Rewriting, events, and causal order}

\begin{definition}[Non-deleting rewriting rule]
A \emph{rule} is a monomorphism $L \hookrightarrow R$ of finite simplicial
sets. An \emph{application} of the rule to a configuration $X$ is a choice of
match $m : L \to X$; the result is the pushout $X \cup_L R$. Because $L
\hookrightarrow R$ is a monomorphism and nothing is removed, applications
only add cells. (This is the non-deleting special case of double-pushout
graph rewriting \citep{ehrig2006dpo}, where the interface equals the
left-hand side.)
\end{definition}

\begin{definition}[Run; event; causal order]
\label{def:causalorder}
A \emph{run} is a finite sequence of rule applications starting from an
initial object. Each application is an \emph{event}. Event $b$ \emph{depends}
on event $a$ (written $a \prec b$) if the match of $b$ uses a cell created by
$a$; the reflexive-transitive closure $\le$ of $\prec$ is the \emph{causal
order}. The \emph{causal past} of an event $e$ is
$\past{e} = \{a : a \le e\}$.
\end{definition}

\begin{remark}
The triple (events, causality, independence) is an event structure in the
sense of \citet{winskel1986event}; we use only the partial order.
\end{remark}

\begin{definition}[Linearization; frame]
A \emph{linearization} (or \emph{linear extension}) of a finite poset
$(P,\le)$ is a total order on $P$ extending $\le$. We write $\Ext(P)$ for the
number of linear extensions. In this framework a \emph{reference frame} is a
linearization of the causal order.
\end{definition}

\subsection{Associahedra and the coherence tower}

\begin{definition}[Catalan numbers; associahedron]
The $n$-th Catalan number is $C_n = \binom{2n}{n}/(n+1)$; $C_{n-1}$ counts
the full bracketings of $n$ composable edges. The \emph{associahedron} $K_n$
\citep{stasheff1963} is the polytope whose vertices are these bracketings
and whose faces encode partial bracketings; $K_4$ is the pentagon.
\end{definition}

The relevance is dynamical: filling $2$-horns creates composites, whose
alternative bracketings pose $3$-horn questions, whose answers pose
compatibility ($4$-horn, pentagon) questions, and so on. Solving promotes
questions up this tower; it never depletes them. Since the number of
bracketings grows like $C_{n-1} \sim 4^n / n^{3/2}$, question supply grows
exponentially with tower level.
